\begin{lemma}[Normalization to a maximum terminal closure]
Let \(G\) be a finite loopless multigraph with individual edge identities,
let \(e_0\) be a distinguished edge, and let the palette have \(k\) colors.
Let \(\mathcal E\) be a positive-rung
\noderef{ExtendedTashkinovTree} with rung number \(n\), displayed tree
\(T\), generating coloring \(\varphi_n\), and certified series
\(\mathcal T\).  Write \(T_n\) for its exact terminal ladder tree,
\(D_n\) for its computed terminal \noderef{TashkinovStoredColors}, and
\(T_{n+1}\) for its exact final series tree.  Suppose
\[
 T=T_{n+1},\qquad \noderef{MaximumProperty}(\mathcal E),\qquad
 \noderef{MaximumDegree}(G)+1\leq k,\qquad
 11\leq\noderef{TashkinovOrder}(G).
\]
Then there exist a partial coloring \(\sigma_n\), tree-sequences
\(T^*\) and \(S\), and a
\noderef{CorrespondingExtendedTashkinovTree} certificate \(\mathcal C\)
from \(\mathcal E\) to the displayed tree \(T^*\) under \(\sigma_n\)
such that all of the following hold:
\begin{enumerate}
\item \(\sigma_n\) is proper, leaves exactly \(e_0\) uncolored, and is
      \((T_n,D_n,\varphi_n)\)-\noderef{StableColoring};
\item exactly one of the following route alternatives is exhibited: the
      terminal kind is PE and \(S=T_n\); or the terminal kind is SE or RE
      and there is an exact edge identity \(f_n\) for which the terminal
      connecting-edge set is \(\{f_n\}\) and
      \noderef{TreeSequenceAdjoinEdge}\((T_n,f_n,S)\) holds;
\item \(T^*\) is a \noderef{TashkinovClosure} of \(S\) under \(\sigma_n\);
\item the exact final tree of the replacement series in \(\mathcal C\)
      is \(T^*\); and
\item the replacement ETT satisfies both \noderef{MaximumProperty} and
      \noderef{TerminalClosureMaximum}.
\end{enumerate}
\end{lemma}
\begin{proof}
Put \(a=\operatorname{Fin.last}(n)\), so the terminal ladder tree,
coloring, and stored set are
\[
 T_n=\mathcal T.\mathrm{tree}(a.\mathrm{castSucc}),\qquad
 \varphi_n=\mathcal T.\mathrm{coloring}(a),\qquad
 D_n=\noderef{TashkinovStoredColors}(\mathcal T,a).
\]
The positivity field of \noderef{ExtendedTashkinovTree} gives \(n\geq1\),
so \(a\) is the successor of the index \(n-1\).  Inspect the certificate
at \(a\), supplied by \noderef{TashkinovSeries} and classified by
\noderef{TashkinovSeriesStep}.  The initial constructor would force
\(n=0\), hence is impossible.  In the SE constructor, its selected
coloring is proper and exactly \(e_0\)-uncolored, its connecting set is
the singleton containing the selected edge, it supplies the exact
one-edge extension \(S=T_n+f_n\), and its closed-extension field says
that \(T_{n+1}\) is a closure of \(S\).  The RE constructor gives the
same four facts with its available revisiting edge.  In the PE constructor,
the recorded current coloring is proper and exactly \(e_0\)-uncolored,
the kind is PE, the seed is exactly \(S=T_n\), and the closed-extension
field says that \(T_{n+1}\) is a closure of this seed.  Thus in every
possible case we have an exact terminal seed \(S\), a proper,
exactly-one-uncolored coloring \(\varphi_n\), the asserted PE or SE/RE
seed route, and
\[
                 \noderef{TashkinovClosure}(\varphi_n,S,T_{n+1}). \tag{1}
\]

By the reflexive clause of \noderef{sc2}, \(\varphi_n\) is
\((T_n,D_n,\varphi_n)\)-stable.  Consequently (1) is a nonempty eligible
family for \noderef{MaximumStableClosureExists}.  Apply that lemma to
choose a proper, exactly-\(e_0\)-uncolored coloring \(\sigma_n\), stable
relative to \((T_n,D_n,\varphi_n)\), and a closure \(T^*\) of the same
exact seed \(S\), such that
\[
 |V(U)|\leq |V(T^*)|                                           \tag{2}
\]
for every eligible coloring \(\rho\) and every closure \(U\) of \(S\)
under \(\rho\).  Applying (2) to \((\varphi_n,T_{n+1})\) from (1) gives
\(|V(T_{n+1})|\leq |V(T^*)|\).  The vertices displayed by a
\noderef{TreeSequence} are injectively indexed, so vertex-set cardinality
is its length plus one; hence
\(|T_{n+1}|\leq |T^*|\).  The strict penultimate containment in
\noderef{ExtendedTashkinovTree}, together with the hypothesis
\(T=T_{n+1}\), gives \(|T_n|<|T_{n+1}|\), and therefore
\[
                             |T_n|<|T^*|.                       \tag{3}
\]

The augmentation part of the selected closure and
\noderef{TashkinovAugmentationSequenceInitialSegment} show that \(S\) is
an exact ordered initial segment of \(T^*\).  In the PE case \(S=T_n\).
In the SE/RE case, the exact one-edge relation makes \(T_n\) an initial
segment of \(S\); composing the two entrywise initial-segment relations
shows that \(T_n\) is an initial segment of \(T^*\).  The same closure
data give precisely the kind-dependent eligibility premise in
\noderef{ModColoring}: direct augmentation from \(T_n\) in PE, or exact
adjoining of the singleton connecting edge followed by augmentation in
SE/RE.  Assertion (vi) of \noderef{TechnicalTheoremLarge}, applied using
the assumed maximum property and numerical bounds, says that
\(\sigma_n\) is a mod coloring.  Its universal clause, applied to the
initial-segment relation, strictness (3), and this exact eligibility route,
therefore supplies a corresponding certificate \(\mathcal C\) on
\(T^*\).

The chosen seed route and selected closure satisfy the hypotheses of
\noderef{CorrespondingTerminalClosureEq}.  That lemma identifies the
exact final output tree of \(\mathcal C\)'s replacement series with
\(T^*\).  Moreover, \noderef{MaximumPropertyTransport}, applied to the
mod-coloring fact, \(\mathcal C\), and the original maximum property,
gives \noderef{MaximumProperty} for the replacement ETT.  It remains to
prove its terminal maximum property.

Write \(\mathcal T^*\) for the replacement series.  The fields of
\noderef{CorrespondingExtendedTashkinovTree} identify its rung with
\(n\), its terminal ladder tree with \(T_n\), its raw stored sets and
terminal extension kind with those of \(\mathcal T\), and its terminal
connecting-edge set with the same exact set of edge identities.  They
also say that its terminal coloring \(\sigma_n^*\) is proper, leaves
exactly \(e_0\) uncolored, and is
\((T_n,D_n,\varphi_n)\)-stable.  The last stability statement identifies
the missing-color set at every vertex of \(T_n\).  Taking the union over
these common vertices identifies the two terminal tree-missing-color
sets.  Since \noderef{TashkinovStoredColors} subtracts exactly that set
from the union of the raw stored sets through the terminal index, the raw
set agreements imply
\[
 D_n^*=D_n.                                                     \tag{4}
\]

Now take an arbitrary competitor \(\rho\) in the definition of
\noderef{TerminalClosureMaximum} for the replacement ETT, with prescribed
seed \(S'\) and closure \(U\).  Using terminal ladder equality and (4),
its assumed \((T_n^*,D_n^*,\sigma_n^*)\)-stability is stability relative
to \((T_n,D_n,\sigma_n^*)\).  Compose it, in this order, after the
certificate's \((T_n,D_n,\varphi_n)\)-stability.  The transitivity clause
of \noderef{sc2} yields that \(\rho\) is
\((T_n,D_n,\varphi_n)\)-stable.                                \tag{5}

We next identify \(S'\) with the selected seed \(S\).  Kind agreement
rules out every mismatch between PE and SE/RE.  In the matching PE case,
both seeds are literally the common terminal ladder tree.  In the matching
SE/RE case, equality of the two singleton connecting-edge sets identifies
their exact edge identities.  Both \(S'\) and \(S\) are therefore obtained
by adjoining the same edge identity to the same ordered tree, so
\noderef{TreeSequenceAdjoinEdgeUnique} gives \(S'=S\).  Thus the competitor
closure is a closure of \(S\), and (2), using (5) together with its
properness and exact-one-uncoloredness, gives
\(|V(U)|\leq |V(T^*)|\).  Finally,
\noderef{CorrespondingTerminalClosureEq} already identified \(T^*\) with
the replacement series' exact final tree.  Rewriting the right side by
that equality proves the required universal bound in
\noderef{TerminalClosureMaximum}.  All five asserted conclusions now hold.
\end{proof}
